% GENERATED FILE. Edit canonical lessons, then run npm run generate:books.
% canonical-source-hash: fnv1a64:ae1df171145721fc
% canonical-lessons: PT-C419-cruzeiro, PT-C419-camarote, PT-C419-marinheiro, PT-C419-proa, PT-C419-popa

\chapter{No navio}
\label{ch:pt-no-navio}
\hlchaptermodality{419}

\begin{chapteropening}
\textbf{By the end of this chapter:} \emph{I can talk about a voyage: the cruise, the cabin, the sailor, the bow and the stern.}

Complete the last lesson of chapter 419: I can talk about a voyage: the cruise, the cabin, the sailor, the bow and the stern.
\end{chapteropening}

\section[o cruzeiro]{o cruzeiro — the cruise; the stone cross (by a road)}

\label{lesson:PT-C419-cruzeiro}



\begin{quote}
Before the new one: say the Portuguese for the canoe, then the Portuguese for the boatman.
\end{quote}

\subsection*{You'll want to know: o cruzeiro}
\textbf{o cruzeiro} — \textquotedblleft{}the cruise; the stone cross (by a road)\textquotedblright{}.

\begin{quote}
Os meus pais fizeram um cruzeiro de uma semana.
\end{quote}

\begin{quote}
\emph{My parents went on a week-long cruise.}
\end{quote}

In villages, \emph{um cruzeiro} is a tall stone cross at a crossroads.

\begin{cousinweb}
\emph{cruzeiro} is built on \emph{cruz} 'cross', inherited from Latin \emph{crux, crucis} (Lewis \& Short); EtymDB lists Spanish \emph{crucero} beside it. A cruise crosses the sea back and forth.
\end{cousinweb}

\begin{grammarlens}[title={What you've built}]
One more word for telling, explaining and describing at B1.
\end{grammarlens}

\subsection*{Your turn}
\begin{itemize}
\raggedright
  \item \emph{Say it:} \emph{o cruzeiro}
  \item \emph{Say it:} \emph{o cruzeiro}, once more
  \item \emph{Say it:} a sentence of your own with \emph{o cruzeiro}
  \item \emph{Recall:} say the Portuguese for the canoe, then the Portuguese for the boatman, then say \emph{o cruzeiro} again
\end{itemize}

\begin{tcolorbox}[breakable,colback=teal!4,colframe=teal!35!black,title={Before you move on}]
What does o cruzeiro mean? (\textquotedblleft{}The cruise; the stone cross (by a road)\textquotedblright{}.) Say it once more.
\end{tcolorbox}
\section[o camarote]{o camarote — the cabin (on a ship); the box (in a theatre)}

\label{lesson:PT-C419-camarote}



\begin{quote}
Before the new one: say the Portuguese for the boatman, then the Portuguese for the cruise.
\end{quote}

\subsection*{You'll want to know: o camarote}
\textbf{o camarote} — \textquotedblleft{}the cabin (on a ship); the box (in a theatre)\textquotedblright{}.

\begin{quote}
O nosso camarote tem uma janela redonda.
\end{quote}

\begin{quote}
\emph{Our cabin has a round window.}
\end{quote}

At a stadium or a concert, \emph{o camarote} is the VIP box.

\begin{cousinweb}
\emph{camarote} is built on \emph{câmara} 'chamber' with the small-size ending \emph{-ote} (EtymDB); \emph{câmara} goes back to Latin \emph{camera}, from Greek \emph{kamára} 'vaulted room'.
\end{cousinweb}

\begin{grammarlens}[title={What you've built}]
One more word for telling, explaining and describing at B1.
\end{grammarlens}

\subsection*{Your turn}
\begin{itemize}
\raggedright
  \item \emph{Say it:} \emph{o camarote}
  \item \emph{Say it:} \emph{o camarote}, once more
  \item \emph{Say it:} a sentence of your own with \emph{o camarote}
  \item \emph{Recall:} say the Portuguese for the boatman, then the Portuguese for the cruise, then say \emph{o camarote} again
\end{itemize}

\begin{tcolorbox}[breakable,colback=teal!4,colframe=teal!35!black,title={Before you move on}]
What does o camarote mean? (\textquotedblleft{}The cabin (on a ship); the box (in a theatre)\textquotedblright{}.) Say it once more.
\end{tcolorbox}
\section[o marinheiro]{o marinheiro — the sailor}

\label{lesson:PT-C419-marinheiro}



\begin{quote}
Before the new one: say the Portuguese for the cruise, then the Portuguese for the cabin.
\end{quote}

\subsection*{You'll want to know: o marinheiro}
\textbf{o marinheiro} — \textquotedblleft{}the sailor\textquotedblright{}.

\begin{quote}
Um marinheiro ajudou-nos com as malas.
\end{quote}

\begin{quote}
\emph{A sailor helped us with our suitcases.}
\end{quote}

\emph{Marinheiro de primeira viagem} — a beginner.

\begin{cousinweb}
\emph{marinheiro} is built on \emph{marinha} 'seashore, navy', from Latin \emph{marīnus} 'of the sea', on \emph{mare} 'sea' (Lewis \& Short). EtymDB lists French \emph{marinier} and Italian \emph{marinaio} as cousins.
\end{cousinweb}

\begin{grammarlens}[title={What you've built}]
One more word for telling, explaining and describing at B1.
\end{grammarlens}

\subsection*{Your turn}
\begin{itemize}
\raggedright
  \item \emph{Say it:} \emph{o marinheiro}
  \item \emph{Say it:} \emph{o marinheiro}, once more
  \item \emph{Say it:} a sentence of your own with \emph{o marinheiro}
  \item \emph{Recall:} say the Portuguese for the cruise, then the Portuguese for the cabin, then say \emph{o marinheiro} again
\end{itemize}

\begin{tcolorbox}[breakable,colback=teal!4,colframe=teal!35!black,title={Before you move on}]
What does o marinheiro mean? (\textquotedblleft{}The sailor\textquotedblright{}.) Say it once more.
\end{tcolorbox}
\section[a proa]{a proa — the bow (of a ship)}

\label{lesson:PT-C419-proa}



\begin{quote}
Before the new one: say the Portuguese for the cabin, then the Portuguese for the sailor.
\end{quote}

\subsection*{You'll want to know: a proa}
\textbf{a proa} — \textquotedblleft{}the bow (of a ship)\textquotedblright{}.

\begin{quote}
Da proa, vemos as ilhas ao longe.
\end{quote}

\begin{quote}
\emph{From the bow, we can see islands in the distance.}
\end{quote}

\emph{À proa} — at the front of the boat.

\begin{cousinweb}
\emph{proa} comes from Latin \emph{prōra} 'bow of a ship', from Greek \emph{prōira} (EtymDB).
\end{cousinweb}

\begin{grammarlens}[title={What you've built}]
One more word for telling, explaining and describing at B1.
\end{grammarlens}

\subsection*{Your turn}
\begin{itemize}
\raggedright
  \item \emph{Say it:} \emph{a proa}
  \item \emph{Say it:} \emph{a proa}, once more
  \item \emph{Say it:} a sentence of your own with \emph{a proa}
  \item \emph{Recall:} say the Portuguese for the cabin, then the Portuguese for the sailor, then say \emph{a proa} again
\end{itemize}

\begin{tcolorbox}[breakable,colback=teal!4,colframe=teal!35!black,title={Before you move on}]
What does a proa mean? (\textquotedblleft{}The bow (of a ship)\textquotedblright{}.) Say it once more.
\end{tcolorbox}
\section[a popa]{a popa — the stern (of a ship)}

\label{lesson:PT-C419-popa}



\begin{quote}
Before the new one: say the Portuguese for the sailor, then the Portuguese for the bow.
\end{quote}

\subsection*{You'll want to know: a popa}
\textbf{a popa} — \textquotedblleft{}the stern (of a ship)\textquotedblright{}.

\begin{quote}
As crianças ficaram na popa a ver as ondas.
\end{quote}

\begin{quote}
\emph{The children stayed at the stern watching the waves.}
\end{quote}

\emph{De vento em popa} — going very well, with the wind behind you.

\begin{cousinweb}
\emph{popa} is inherited from Latin \emph{puppis} 'stern of a ship' (Lewis \& Short).
\end{cousinweb}

\begin{grammarlens}[title={What you've built}]
One more word for telling, explaining and describing at B1.
\end{grammarlens}

\subsection*{Your turn}
\begin{itemize}
\raggedright
  \item \emph{Say it:} \emph{a popa}
  \item \emph{Say it:} \emph{a popa}, once more
  \item \emph{Say it:} a sentence of your own with \emph{a popa}
  \item \emph{Recall:} say the Portuguese for the sailor, then the Portuguese for the bow, then say \emph{a popa} again
\end{itemize}

\begin{tcolorbox}[breakable,colback=teal!4,colframe=teal!35!black,title={Before you move on}]
What does a popa mean? (\textquotedblleft{}The stern (of a ship)\textquotedblright{}.) Say it once more.
\end{tcolorbox}
